\documentclass[12pt]{article}
\usepackage[utf8]{inputenc}
\usepackage[T1]{fontenc}
\usepackage{hyperref}
\title{V\~oro Language Resources for K-12 Curriculum Development}
\author{Compiled by Claude Sonnet 5 (claude-sonnet-5) for Andr\'as Kornai}
\date{}
\begin{document}
\maketitle
\section{Introduction}
V\~oro (ISO 639-3 vro; Glottocode sout2679) is a Uralic language of south-eastern Estonia, with about 100,000 speakers. Unlike many languages in this survey, it has an established institutional base in the V\~oro Institute and some school instruction. The search was brief, and the Institute's own catalogue (wi.ee) was not among the results, so the list below is not exhaustive.
\section{Primary Educational Resources}
\begin{itemize}
\item A new V\~oro-language children's textbook with 30 chapters, illustrations by Marja-Liisa Plats and QR codes linking to songs and stories, reported by ERR: \url{https://www.err.ee/1609160458/vorokesed-soovivad-uue-opikuga-voru-keelt-lastele-lahemale-tuua}
\item Curriculum for the 2025 in-service course for V\~oro language-nest (keelepesa) teachers, 108 academic hours: \url{https://wi.ee/wp-content/uploads/2025/01/Voru-keelepesa-opetajate-taiendkoolituse-oppekava-2025-aastal.pdf}
\end{itemize}
\section{Linguistic and Reference Materials}
\begin{itemize}
\item Overview of the V\~oro Institute: \url{https://en.wikipedia.org/wiki/V%C3%B5ro_Institute}
\item Kadri Koreinik and Sulev Iva on V\~oro language status and use: \url{https://keel.ut.ee/sites/default/files/2023-09/09_kadri_koreinik_sulev_iva.pdf}
\end{itemize}
No dictionary or grammar was surfaced by this search.
\section{Religious and Community Resources}
The V\~oro Institute is the main community and state body for the language (see the Wikipedia link above). News of a V\~oro-language children's musical by Voru school pupils: \url{https://www.err.ee/1608245283/voru-koolinoored-toid-lavale-vorukeelse-lastemuusikali}. No scripture resources were found in this search.
\section{Digital and Audio Resources}
Local language-week coverage: \url{https://www.err.ee/1608781483/vorus-on-nadala-valtel-tahelepanu-keskmes-kohalik-keel}. The textbook above uses QR-linked audio. No stand-alone corpus or software was found here.
\section{Curriculum Development Context}
V\~oro teaching in schools started in the mid-1990s, and the state has supported V\~oro language and culture teaching since 2000. A 2018 study found that fewer than half of local schools offered it. The Institute's teacher training for language nests suggests a base for early-years immersion. A fuller survey of the Institute's publications is needed.
\end{document}